\section{Set \& class theory}
Notation:\\

Starting here and for the rest of the document, we will use usual symbol of equality "=" for quasi-equality (before it was "$\simeq$") and "$\equiv$" for equality (before it was "=").

\begin{definition}
\label{def:identity}
Identity.
\[
I = \{(x,x)|\top\}
\]
\end{definition}

\begin{theorem}
\label{thm:identity is relation}
Identity is a relation.
\[
(I \relation)
\]
\[(I\subseteq V\times V)\]
\end{theorem}


\begin{definition}
"F on X".
\[
\mbox{"F on X"} = F\restriction X
\]
\end{definition}


\begin{definition}
\[
(y\in f[A])
\iff
(\exists x\in A: \langle x,y\rangle\in f)
\]
\end{definition}

\begin{definition}
\[
(y\in f\langle A\rangle)
\iff
(\forall x\in A: \langle x,y\rangle\in f)
\]
\end{definition}
\begin{theorem}
\[
X\subseteq \power Y
\iff
\union X \subseteq Y
\]
\end{theorem}
\begin{theorem}
\[
(g\functional)
\land
(M \countable)
\implies
(g[M] \countable)
\]
\end{theorem}
\begin{theorem}
\[
(f^\leftarrow \univocal)
\]
\end{theorem}

\begin{theorem}
\[
(f\injective)
\implies
(f\left[\union M\right] = \union f^{\rightarrow}\left[M\right])
\]
\end{theorem}

\begin{theorem}
\[
(f\univocal)
\implies
f^{-1}\left[\union M\right] =\union f^{\leftarrow}\left[M\right]
\]
\end{theorem}
